\chapter{Classes and Objects}
% This session refactors your session 5 Iris classifier using Object-Oriented Programming (OOP).
% The goal is to bundle the configuration (threshold) and the behavior (prediction logic) together in a class named \texttt{IrisRuleClassifier}.

\section{Introduction}
In session 4 and session 5, you built an Iris classifier using \textbf{functions}.
Your pipeline worked, but you may have noticed a design pattern:
many functions needed the same configuration values (for example, the threshold and the label names), so you passed a dictionary called \texttt{settings} around repeatedly.

\noindent
\textbf{Why refactor?}
That function-based design is fine for one small classifier, but it becomes awkward as soon as you want to run experiments:
\begin{itemize}\setlength{\itemsep}{2pt}
    \item If you want 10 classifiers with 10 thresholds, you now need 10 different \texttt{settings} dictionaries.
    \item Every time you call a function, you must remember to pass the correct \texttt{settings}.
    \item If you add one new setting later, you must update many function signatures and many function calls.
\end{itemize}

\noindent
This session (Session 6), you refactor your session 5 design using \textbf{Object-Oriented Programming (OOP)}.
The core idea is simple:
\begin{quote}
\textbf{Bundle the data (configuration) and the behavior (functions) together.}
\end{quote}

\noindent
In Python, that bundle is called a \textbf{class}.
\begin{itemize}\setlength{\itemsep}{2pt}
    \item A \textbf{class} is a blueprint that describes what data and behavior a thing should have.
    \item An \textbf{object} is a single instance created from that blueprint.
    \item An \textbf{attribute} is data stored on the object (for example, \texttt{self.threshold}).
    \item A \textbf{method} is a function stored on the class that operates on the object (for example, \texttt{predict()} and \texttt{evaluate()}).
\end{itemize}

\noindent
\textbf{Refactoring promise:} the classifier rule itself stays the same.
This session reorganizes code, but it should \textbf{not change} what ``setosa'' means or how the threshold rule works.

\noindent
\textbf{The Goal:} refactor your session 5 classifier into a class named \texttt{IrisRuleClassifier} so each object ``remembers'' its own threshold.

\subsection{Transition Map: session 5 Functions $\rightarrow$ session 6 Class}
The classifier logic itself (``if petal\_length $<$ threshold \dots'') remains exactly the same.
What changes is \textit{where} that logic lives and \textit{how} you access it.

\noindent
Use this table as your translation guide. It shows how session 5 functions and settings translate into session 6 attributes and methods:

\input{session6/tab6.tex}

\subsection{Deliverable}
Open or create \texttt{session6/session6\_iris\_template.py}.
By the end of this session, it will contain:
\begin{itemize}\setlength{\itemsep}{2pt}
    \item a complete \texttt{IrisRuleClassifier} class (constructor + methods),
    \item a small \texttt{main()} script that creates objects and evaluates them.
\end{itemize}
\noindent
Think of this as turning your session 5 pipeline into a reusable component you can instantiate many times.


\section{The Blueprint and Data (\texttt{\_\_init\_\_})}

\subsection{Step 1: The Blueprint and Data (\texttt{\_\_init\_\_})}
A class starts with a special constructor method called \texttt{\_\_init\_\_}.
This method runs \textbf{automatically} when you create a new object.
Its job is to put the object into a valid starting state.

\noindent
In session 5, you used \texttt{get\_default\_settings()} to build a dictionary.
In session 6, you store that information as \textbf{attributes} on \texttt{self}.
\texttt{self} is how a method refers to the current object.

\begin{minted}[fontsize=\small, linenos]{python}
class IrisRuleClassifier:
    def __init__(self, threshold=2.0):
        # Store configuration inside the instance (self)
        self.threshold = threshold
        
        # Store labels so they are not hardcoded in multiple places
        self.positive_label = "setosa"
        self.negative_label = "not_setosa"
\end{minted}

\noindent
\textbf{Task:}
\begin{enumerate}
    \item Define this class in your file.
    \item Below the class definition, create an instance: \texttt{my\_clf = IrisRuleClassifier(threshold=2.5)}.
    \item Print the stored threshold: \texttt{print("Threshold:", my\_clf.threshold)}.
\end{enumerate}
\noindent
\textbf{Before you code (reasoning):} list the values that were previously inside \texttt{settings}.
Those values are good candidates to become attributes on the class.

\noindent
\textbf{After you code (verification):} you should be able to create two objects with different thresholds and see that each object stores its own value.
\begin{minted}[fontsize=\small, linenos]{python}
clf_a = IrisRuleClassifier(threshold=1.5)
clf_b = IrisRuleClassifier(threshold=2.5)
print(clf_a.threshold, clf_b.threshold)  # should print two different numbers
\end{minted}
\noindent
\textbf{Common beginner pitfall:} forgetting \texttt{self}.
Inside methods you must write \texttt{self.threshold}, not just \texttt{threshold}, because \texttt{threshold} is not a global variable.


\section{Method in (\texttt{class})}

\subsection{Step 2: The Logic (\texttt{predict})}
In session 5, you wrote \texttt{classify\_sample(sample, settings)}.
In session 6, this becomes a \textbf{method}.
Because \texttt{threshold} and the labels already live on \texttt{self}, you no longer need a \texttt{settings} dictionary argument.

\noindent
Conceptually, you are changing:
\begin{quote}
\texttt{classify\_sample(sample, settings)} $\rightarrow$ \texttt{my\_clf.predict(sample)}
\end{quote}
Instead of passing configuration into every call, you ask the configured object to do the work.

\begin{minted}[fontsize=\small, linenos]{python}
    def predict(self, sample):
        # We assume sample is a dict with 'petal_length'
        if sample["petal_length"] < self.threshold:
            return self.positive_label
        return self.negative_label
\end{minted}

\noindent
\textbf{Task:} Add this method to your class (indented inside the class).
\noindent
\textbf{Before you code (reasoning):} identify which inputs are ``per sample'' and which inputs are ``settings''.
\texttt{sample} changes every iteration, so it stays an argument. \texttt{threshold} is a setting, so it becomes \texttt{self.threshold}.

\noindent
\textbf{After you code (verification):} call \texttt{my\_clf.predict(...)} on a single dictionary and make sure you get either \texttt{"setosa"} or \texttt{"not\_setosa"}.

\subsection{Step 3: Processing the Dataset (Evaluation)}
In session 5, you had a function \texttt{initialize\_predictions} that looped over the data.
We will now create a method \texttt{evaluate(self, dataset)} that does the entire prediction loop and calculates accuracy.

\noindent
\textbf{Why make this a method?}
Because evaluation depends on \texttt{predict()}, and \texttt{predict()} depends on the object's configuration.
By keeping them together, you avoid passing settings around and you reduce the chance of mixing up configurations.

\begin{minted}[fontsize=\small, linenos]{python}
    def evaluate(self, dataset):
        correct = 0
        total = 0
        
        for sample in dataset:
            # 1. Predict using the class method
            prediction = self.predict(sample)
            
            # 2. Determine Truth (Simplified for this exercise)
            # Truth is 'setosa' if the species is 'setosa'
            truth = "setosa" if sample["species"] == "setosa" else "not_setosa"
            
            # 3. Score
            if prediction == truth:
                correct += 1
            total += 1
            
        return correct / total if total > 0 else 0.0
\end{minted}

\noindent
\textbf{Before you code (reasoning):}
\begin{itemize}\setlength{\itemsep}{2pt}
    \item The evaluation loop should not re-implement prediction logic. It should call \texttt{self.predict(sample)}.
    \item The method should keep counting logic local (\texttt{correct}, \texttt{total}).
    \item The method should return one clear metric (accuracy), so it is easy to compare classifiers.
\end{itemize}

\noindent
\textbf{After you code (verification):}
\begin{itemize}\setlength{\itemsep}{2pt}
    \item Accuracy is a number between 0.0 and 1.0.
    \item If you change the object's threshold, accuracy can change without modifying the evaluation code.
    \item If the dataset is empty, you should not crash (avoid division by zero).
\end{itemize}

\noindent
\textbf{Common beginner pitfall:} forgetting to increment \texttt{total} for every sample, even wrong ones.


\section{Still not sure the best section name for this}


\subsection{Step 4: Putting it all together (Main Script)}
Now you can use your class to define a small dataset and test different classifiers easily.
The workflow becomes:
\begin{enumerate}
    \item Create a dataset (list of dictionaries).
    \item Create one or more classifier objects (each with its own threshold).
    \item Ask each object to evaluate itself on the dataset.
\end{enumerate}

\begin{minted}[fontsize=\small, linenos]{python}
# 1. Define dataset (a small test set)
dataset = [
    {"petal_length": 1.4, "species": "setosa"},
    {"petal_length": 1.5, "species": "setosa"},
    {"petal_length": 4.5, "species": "versicolor"},
    {"petal_length": 6.0, "species": "virginica"},
]

# 2. Create Objects with different configurations
clf_strict = IrisRuleClassifier(threshold=1.5)   # Strict threshold
clf_loose  = IrisRuleClassifier(threshold=2.5)   # Loose threshold

# 3. Evaluate each one
print("Strict Classifier Accuracy:", clf_strict.evaluate(dataset))
print("Loose Classifier Accuracy: ", clf_loose.evaluate(dataset))
\end{minted}
\noindent
\textbf{What to notice:} the dataset is passed in, but the threshold is not.
The threshold lives inside each object.

\section{Why is this better than session 5?}
In session 5, if we wanted 2 different thresholds, we had to manage 2 different \texttt{settings} dictionaries and ensure we passed the correct one to every function call.
In session 6, each object \textbf{remembers its own settings}.
We simply ask the object to \texttt{evaluate(dataset)}.

\noindent
This design scales better because:
\begin{itemize}\setlength{\itemsep}{2pt}
    \item You can create many objects with different configurations and compare them cleanly.
    \item You reduce the number of parameters passed around between functions.
    \item You get a natural place to store state and behavior together, which prepares you for Session 7 (where you will create multiple classifier types).
\end{itemize}

\section{Submission Requirement}
Add your student label at the end:
\begin{minted}[fontsize=\small, linenos]{python}
print("=== Student Label ===")
print("ID: ...")
\end{minted}
Run the script and verify strict vs.\ loose accuracy scores are printed.
\noindent
\textbf{Minimum check:} your output should show at least two different accuracy numbers when you change thresholds, proving that the object stores configuration correctly.

